% © 2026 Ing. David Jaroš — CC BY-NC-ND 4.0
%
% File: research_tracks/T3_ALPHA/em_modular_coupling_identification.tex
% Purpose: Close the physical identification n = alpha^{-1} by matching the
%          UBT electromagnetic winding modulus to the standard compact-U(1)
%          Maxwell coupling modulus.
% Status: [L1 conditional] — closes the final alpha interpretation gap provided
%         the UBT EM projection is accepted as the compact U(1)_EM line bundle.
% Date: 2026-06-13

% UBT-AI-PROVENANCE-BEGIN
% schema: ubt-ai-provenance/v1
% tier: C_working
% ai_assistance: disclosed
% human_review: risk-based
% editorial_responsibility: Ing. David Jaroš
% policy: ../../AI_PROVENANCE.md
% notice: Working material; exhaustive human review is not claimed.
% UBT-AI-PROVENANCE-END

\ifdefined\INCLUDEMODE\else
\documentclass[12pt,a4paper]{article}
\usepackage[a4paper,margin=2.5cm]{geometry}
\usepackage{amsmath,amssymb,amsthm,mathtools}
\usepackage{booktabs}
\usepackage{xcolor}
\usepackage[most]{tcolorbox}
\usepackage{hyperref}

\newtheorem{definition}{Definition}[section]
\newtheorem{lemma}[definition]{Lemma}
\newtheorem{theorem}[definition]{Theorem}
\newtheorem{corollary}[definition]{Corollary}
\newtheorem{remark}[definition]{Remark}

\newtcolorbox{statusbox}[1][]{
  colback=yellow!8!white,
  colframe=orange!70!black,
  arc=3pt,
  boxrule=1pt,
  title=Status,
  #1
}
\newtcolorbox{resultbox}[1][]{
  colback=green!8!white,
  colframe=green!50!black,
  arc=3pt,
  boxrule=1pt,
  title=Result,
  #1
}

\newcommand{\Lone}{\textbf{[L1 cond.]}}
\newcommand{\Std}{\textbf{[STD]}}

\title{Electromagnetic Modular Coupling and the Identification\\
       \large $n = \alpha^{-1}$ in the UBT Alpha Track}
\author{Ing.\ David Jaro\v{s}}
\date{2026-06-13}

\begin{document}
\maketitle
\fi

\begin{abstract}
The numerical UBT alpha fixed point produces a stationary winding stiffness
\(n_*=137.035999177548864\ldots\).  The remaining physical gap is to identify
this dimensionless winding stiffness with the inverse fine-structure constant.
This note shows that, once the UBT electromagnetic projection is treated as the
compact unit-charge \(U(1)_{\rm EM}\) line bundle, the identification follows
from the standard complexified Maxwell coupling
\[
  \tau_{\rm EM}=\frac{\theta_{\rm EM}}{2\pi}+i\frac{4\pi}{e^2}.
\]
The UBT winding modulus is the same compact-U(1) modulus,
\(\tau^{\rm UBT}_{\rm EM}=\chi+i n\).  In the CP-even vacuum \(\chi=0\), hence
\[
  n = \frac{4\pi}{e^2}=\alpha^{-1}.
\]
The result does not change the numerical alpha fixed point.  It only closes the
normalisation/interpretation bridge from the UBT winding index to the physical
low-energy electromagnetic coupling.
\end{abstract}

\begin{statusbox}
\textbf{Classification.} This is a conditional closure theorem.  It is not a
new numerical fit.  The only physical assumption is the standard one: the UBT
EM projection defines the compact unit-charge \(U(1)_{\rm EM}\) line bundle seen
by low-energy electromagnetism.  Under this assumption the normalisation
\(n=\alpha^{-1}\) is fixed by compact-U(1) Maxwell theory and Dirac/Wilson-line
charge quantisation.
\end{statusbox}

\section{Standard compact-U(1) normalisation}

Let \(A\) be a compact unit-charge \(U(1)\) connection on a line bundle \(L\),
with curvature \(F=dA\) normalised by
\[
  \frac{1}{2\pi}\int_{\Sigma}F \in \mathbb Z
\]
for every closed two-cycle \(\Sigma\).  In Euclidean signature the Maxwell
action is
\begin{equation}
  S_{\rm Max}[A]
  =
  \frac{1}{2e^2}\int F\wedge *F
  - i\frac{\theta_{\rm EM}}{8\pi^2}\int F\wedge F .
  \label{eq:stdMaxwell}
\end{equation}
Equivalently, in tensor notation,
\[
  \frac{1}{2e^2}\int F\wedge *F
  =
  \frac{1}{4e^2}\int F_{\mu\nu}F^{\mu\nu}\,d^4x .
\]
The standard complexified coupling is
\begin{equation}
  \tau_{\rm EM}
  =
  \frac{\theta_{\rm EM}}{2\pi}
  + i\frac{4\pi}{e^2}.
  \label{eq:tauEM}
\end{equation}
Since \(\alpha=e^2/(4\pi)\),
\begin{equation}
  {\rm Im}\,\tau_{\rm EM}=\frac{4\pi}{e^2}=\alpha^{-1}.
  \label{eq:imalpha}
\end{equation}
This is a normalisation identity, not a phenomenological input.

\section{UBT electromagnetic winding modulus}

The UBT alpha track defines a compact \(\psi\)-phase sector.  Its EM projection
selects the unit-charge Abelian phase bundle that survives at low energy as
\(U(1)_{\rm EM}\).  The winding stiffness \(n\) is the coefficient of the
CP-even Hodge norm of this compact connection.  The possible CP-odd phase is
written \(\chi\).  Thus the UBT EM sector has the compact-U(1) action
\begin{equation}
  S^{\rm UBT}_{\rm EM}[A]
  =
  \frac{n}{8\pi}\int F\wedge *F
  - i\frac{\chi}{4\pi}\int F\wedge F .
  \label{eq:UBTMaxwell}
\end{equation}
In tensor notation,
\begin{equation}
  S^{\rm UBT}_{\rm EM}[A]
  =
  \frac{n}{16\pi}\int F_{\mu\nu}F^{\mu\nu}\,d^4x
  - i\frac{\chi}{4\pi}\int F\wedge F .
  \label{eq:UBTMaxwellComponents}
\end{equation}
Equation~\eqref{eq:UBTMaxwell} is the compact-U(1) Hodge norm with dimensionless
modulus
\begin{equation}
  \tau^{\rm UBT}_{\rm EM}=\chi+i n .
  \label{eq:tauUBT}
\end{equation}
The same \(n\) is the variable whose vacuum value is selected by the eta/Layer2
fixed-point equation, because the fixed point extremises the partition function
of this compact \(\psi\)-winding sector.

\section{Identification theorem}

\begin{theorem}[Electromagnetic modular identification \Lone]
Assume that the UBT electromagnetic projection is the compact unit-charge
\(U(1)_{\rm EM}\) line bundle of the low-energy theory, and that the stationary
UBT winding stiffness \(n\) is the imaginary part of its compact-U(1) modulus:
\[
  \tau^{\rm UBT}_{\rm EM}=\chi+i n.
\]
Then, after matching to the standard Maxwell coupling,
\[
  n=\frac{4\pi}{e^2}=\alpha^{-1}.
\]
\end{theorem}

\begin{proof}
Compare the CP-even term in the UBT action~\eqref{eq:UBTMaxwell} with the
standard Maxwell action~\eqref{eq:stdMaxwell}:
\[
  \frac{n}{8\pi}\int F\wedge *F
  =
  \frac{1}{2e^2}\int F\wedge *F .
\]
Therefore
\[
  \frac{n}{8\pi}=\frac{1}{2e^2},
  \qquad
  n=\frac{4\pi}{e^2}.
\]
Using \(\alpha=e^2/(4\pi)\), we obtain
\[
  n=\alpha^{-1}.
\]
The CP-odd terms also match if \(\chi=\theta_{\rm EM}/(2\pi)\):
\[
  -i\frac{\chi}{4\pi}\int F\wedge F
  =
  -i\frac{\theta_{\rm EM}}{8\pi^2}\int F\wedge F .
\]
Thus the full complexified coupling agrees:
\[
  \tau^{\rm UBT}_{\rm EM}=\tau_{\rm EM}.
\]
In the CP-even low-energy vacuum \(\theta_{\rm EM}=0\), hence \(\chi=0\) and
\(\tau^{\rm UBT}_{\rm EM}=i n=i\alpha^{-1}\).
\end{proof}

\begin{corollary}[Closure of the final alpha interpretation gap \Lone]
If the eta/Layer2 fixed point gives
\[
  n_*=137.035999177548864\ldots,
\]
then the low-energy electromagnetic coupling predicted by UBT is
\[
  \alpha^{-1}_{\rm UBT}=n_*=137.035999177548864\ldots .
\]
\end{corollary}

\section{Why this is not a fitted normalisation}

The factor \(4\pi\) does not come from tuning.  It is fixed by the standard
compact-U(1) definition of the complexified Maxwell coupling and by the
normalisation of electric charge:
\[
  \tau_{\rm EM}=\frac{\theta}{2\pi}+i\frac{4\pi}{e^2}.
\]
Changing this factor would change the quantisation convention of the physical
unit-charge Wilson line and would no longer describe the observed
\(U(1)_{\rm EM}\) coupling.

\section{Remaining structural condition}

The compact unit-charge EM projection is now reduced to the electroweak-vacuum
condition.  The remaining conditional statement is structural, not numerical:
\[
  \boxed{
  V(\Theta)\text{ has a non-zero minimum in the minimal electroweak doublet sector.}
  }
\]
The \(\Theta_0\) vacuum-orbit theorem then gives the standard
\((SU(2)_L,Y=1)\) representative, the EM projection theorem gives primitive
compact \(U(1)_{\rm EM}\), and this note gives \(n=\alpha^{-1}\) by Maxwell
modular normalisation.

\begin{resultbox}
\textbf{Conclusion.} The alpha track no longer needs a heuristic statement
``let \(n=\alpha^{-1}\).''  It can state the compact-U(1) matching theorem:
\[
  \tau^{\rm UBT}_{\rm EM}=\chi+i n
  \quad\Longleftrightarrow\quad
  \tau_{\rm EM}=\frac{\theta}{2\pi}+i\alpha^{-1}.
\]
At \(\theta=0\), this gives \(n=\alpha^{-1}\).
\end{resultbox}

\ifdefined\INCLUDEMODE\else
\end{document}
\fi
